\chapter{Batteries}

\section{Introduction: Energy Storage and Electrochemistry}

The growing demand for portable energy, electric vehicles, and renewable energy integration has made efficient energy storage a central challenge of our time. Raw energy — whether from fossil fuels, nuclear reactions, solar irradiation, or wind — must often be transformed and stored before use. Among storage technologies, electrochemical batteries offer a unique combination of portability, scalability, and high energy conversion efficiency.

The classical thermomechanical route to electricity — combustion, thermal engine, alternator — is fundamentally limited by the Carnot principle, which caps efficiency at $\eta_{\max} = 1 - T_c/T_h$. The electrochemical route, by contrast, achieves a direct conversion of chemical energy into electrical work, with a theoretical yield given by:

\begin{equation}
    \eta_{\text{th}} = \frac{\Delta G}{\Delta H}
\end{equation}

For a hydrogen/oxygen fuel cell, this yields $\eta_{\max} \approx 95\%$, far exceeding any heat engine operating between practical temperature limits.

The fundamental unit of an electrochemical energy storage system is the \textbf{electrochemical cell}, composed of two electrodes separated by an electrolyte:

\begin{center}
    \textit{electrode / electrolyte / electrode}
\end{center}

At constant temperature and pressure, the Gibbs free energy of the cell reaction is directly related to the cell voltage:

\begin{equation}
    \Delta G(T,P) = -nF\Delta\varepsilon
\end{equation}

where $n$ is the number of electrons exchanged, $F = 96{,}485\,\text{C/mol}$ is Faraday's constant, and $\Delta\varepsilon$ is the electromotive force. The spontaneity condition requires $\Delta\varepsilon > 0$. The maximum achievable cell voltage is:

\begin{equation}
    \Delta V_{\max} = \Delta\varepsilon = -\frac{\Delta G}{nF}
\end{equation}

For example, the lead-acid cell gives $\Delta V_{\max} = 2.1\,\text{V}$ per cell, so a standard 12\,V car battery consists of six cells in series.

\textbf{Rechargeable} (secondary) batteries must additionally satisfy three conditions: the cell reaction must be reversible, the charge/discharge cycling yield must be acceptable, and material degradation must remain within tolerable bounds.

The history of rechargeable batteries spans nearly two centuries, from Planté's lead-acid cell (1859) through nickel-cadmium and nickel-metal hydride technologies, culminating in the lithium-ion battery developed between 1970 and 1990 by Whittingham, Goodenough, and Yoshino — work recognised by the 2019 Nobel Prize in Chemistry. Over this period, the specific energy of rechargeable batteries has risen from roughly 30\,Wh/kg (lead-acid) to over 180\,Wh/kg (Li-ion), with volumetric energy density approaching the theoretical limit of around 400\,Wh/L.

\section{How a Li-ion Battery Works}

\subsection{General Operating Principle}

A lithium-ion battery operates through the \textbf{reversible intercalation} of lithium ions (Li$^+$) between two host electrode materials. Crucially, lithium is never reduced to its metallic form during normal operation; it remains ionic throughout the charge/discharge cycle.

\begin{itemize}
    \item \textbf{Negative electrode (anode):} graphite, which hosts Li$^+$ ions in its layered structure to form LiC$_6$.
    \item \textbf{Positive electrode (cathode):} a lithiated transition metal oxide, typically of the form LiMO$_x$ (e.g., LiCoO$_2$, LiFePO$_4$, LiMn$_2$O$_4$, or NMC variants).
\end{itemize}

The overall cell reaction can be written as:

\begin{equation}
    \text{LiMO}_x + \text{C}_n \underset{\text{charge}}{\stackrel{\text{discharge}}{\rightleftharpoons}} \text{Li}_{1-y}\text{MO}_x + \text{Li}_y\text{C}_n
\end{equation}

During \textbf{discharge}, Li$^+$ ions de-intercalate from the graphite anode, travel through the electrolyte, and re-intercalate into the cathode oxide, while electrons flow through the external circuit producing electrical work. During \textbf{charge}, the process is reversed.

\subsection{Criteria for Electrode Materials}

An effective intercalation electrode material must satisfy several criteria simultaneously:

\begin{enumerate}
    \item \textbf{High Li$^+$ insertion capacity:} the material must accommodate a large number of lithium ions per formula unit, yielding high specific capacity (mAh/g).
    \item \textbf{Reversibility:} insertion and de-insertion must be highly reversible with minimal hysteresis, ensuring long cycle life.
    \item \textbf{Voltage plateau:} the electrode potential should be stable (flat plateau) at an adequate value with respect to Li$^+$/Li.
    \item \textbf{Good electronic and ionic conductivity:} both electron and ion transport through the electrode must be fast enough to support practical current densities.
    \item \textbf{Low volume change:} the crystal structure should expand and contract minimally upon Li$^+$ insertion/extraction to avoid mechanical degradation.
    \item \textbf{Durability and low cost:} the material must be stable over hundreds to thousands of cycles and manufacturable from abundant, inexpensive precursors.
\end{enumerate}

The stored energy of a cell is given by:

\begin{equation}
    E = \Delta U \times i \times t
\end{equation}

where $\Delta U$ is the electromotive force (difference between cathode and anode potentials) and $i \times t$ is the capacity delivered.

\subsection{Representative Electrode Materials}

Table~\ref{tab:electrodes} summarises the key properties of common positive and negative electrode materials.

\begin{table}[ht!]
\centering
\caption{Properties of selected Li-ion electrode materials.}
\label{tab:electrodes}
\begin{tabular}{llll}
\toprule
\textbf{Material} & \textbf{Avg. voltage} & \textbf{Specific capacity} & \textbf{Mass energy} \\
\midrule
\multicolumn{4}{l}{\textit{Positive electrodes}} \\
LiCoO$_2$ & 3.7\,V & 140\,mAh/g & 0.518\,kWh/kg \\
LiMn$_2$O$_4$ & 4.0\,V & 100\,mAh/g & 0.400\,kWh/kg \\
LiFePO$_4$ & 3.3\,V & 150\,mAh/g & 0.495\,kWh/kg \\
\midrule
\multicolumn{4}{l}{\textit{Negative electrodes}} \\
Graphite (LiC$_6$) & 0.1--0.2\,V & 372\,mAh/g & 0.037--0.074\,kWh/kg \\
Titanate (Li$_4$Ti$_5$O$_{12}$) & 1--2\,V & 160\,mAh/g & 0.16--0.32\,kWh/kg \\
Si (Li$_{4.4}$Si) & 0.5--1\,V & 4212\,mAh/g & 2.1--4.2\,kWh/kg \\
\bottomrule
\end{tabular}
\end{table}

Li-ion cells typically deliver 3.5--3.8\,V, compared to 2.0\,V for lead-acid and 1.2\,V for Ni-MH cells. Compared to competing rechargeable technologies, Li-ion offers a superior energy density (180--220\,Wh/kg vs. 30--80\,Wh/kg for older chemistries), faster charging (60--120\,min), lower self-discharge (5--10\%/month), and a charge/discharge round-trip efficiency exceeding 85\%.

\section{Battery Construction}

\subsection{Cell Architecture}

The electrochemical cell consists of thin-film layers stacked or wound together:

\begin{center}
    current collector $\mid$ anode $\mid$ electrolyte $\mid$ separator $\mid$ electrolyte $\mid$ cathode $\mid$ current collector
\end{center}

The anode current collector is copper foil; the cathode current collector is aluminium foil. The separator is a porous polymer membrane that prevents electronic contact between the two electrodes while allowing ionic transport.

\subsection{Composite Electrodes}

Practical electrodes are not monolithic but \textbf{composite structures} deposited as thin films (typically 100\,\textmu m thick) on the metallic current collector. Each electrode layer contains:

\begin{itemize}
    \item \textbf{Active material} (powder): the intercalation compound itself.
    \item \textbf{Binder}: typically polyvinylidene fluoride (PVDF), dissolved in N-methylpyrrolidone (NMP) to form a slurry. PVDF is chemically stable but is a fluorinated polymer and therefore environmentally persistent and difficult to recycle. Alternative water-based binders (cellulose-derived gums, carboxymethyl cellulose) are actively being developed.
    \item \textbf{Conductive additive}: usually carbon black, to ensure electronic percolation through the electrode since many active materials have low intrinsic electronic conductivity.
\end{itemize}

The assembled electrode is soaked with liquid electrolyte, which fills the porous space between particles and ensures ionic access to all active surfaces. Nanostructuring the active material (reducing particle size, introducing porosity) is a key strategy to maximise accessible surface area and reduce solid-state diffusion distances for Li$^+$.

\subsection{Electrolyte}

The electrolyte must be a pure ionic conductor (electron insulator) that is stable across the full voltage window of the cell ($\sim$4\,V for standard Li-ion). Two main types are used:

\begin{itemize}
    \item \textbf{Liquid electrolyte} ("classical" Li-ion): a lithium salt (typically LiPF$_6$) dissolved in a mixture of organic carbonate solvents (e.g., ethylene carbonate, dimethyl carbonate). These electrolytes are highly flammable, requiring the addition of flame-retardant additives and imposing strict purity requirements during manufacture.
    \item \textbf{Gel polymer electrolyte} (Li-ion "polymer"): a polymer matrix swollen with liquid electrolyte, offering improved safety and mechanical flexibility.
\end{itemize}

A critical limitation of liquid electrolytes is that their ionic conductivity decreases significantly at low temperatures, degrading battery performance in cold climates.

\subsection{Cell Formats and Assembly}

Because a single cell provides only 3.5--3.8\,V, practical battery packs assemble many cells in series (to increase voltage) and in parallel (to increase capacity). The cell format can be cylindrical (wound jelly-roll), prismatic (stacked layers), or pouch (flat flexible). Battery manufacturing is a complex multi-step process involving slurry mixing, coating, drying, calendering, slitting, winding or stacking, electrolyte filling, and final sealing, all under tightly controlled atmospheric conditions to exclude moisture and metallic contaminants. Even microscopic metal particles can cause internal short circuits, thermal runaway, and fire.

At the pack level, a \textbf{Battery Management System (BMS)} is essential to monitor cell voltages, temperatures, and state of charge, and to balance cells, prevent overcharge/overdischarge, and manage thermal conditions.

\section{Challenges}

\subsection{Materials Challenges}

The development of Li-ion batteries relies on several materials that are classified as \textbf{critical raw materials} under EU legislation, including lithium, cobalt, nickel, manganese, and natural graphite. These elements are geographically concentrated, subject to supply chain risks, and often associated with environmental and ethical concerns in mining. This drives intensive research into alternative chemistries.

Silicon is a particularly attractive anode candidate, with a theoretical capacity of $\sim$4000\,mAh/g compared to only 372\,mAh/g for graphite. However, Si undergoes a volume expansion of $\sim$320\% upon full lithiation (to Li$_{22}$Si$_5$), causing mechanical cracking and loss of electrical contact upon cycling. Strategies to mitigate this include nanostructured silicon (nanoparticles, nanowires, honeycombs, micropillars), porous composite architectures embedding Si nanoparticles in a compliant carbon matrix, and silicon-graphite blended anodes already entering commercial production.

On the cathode side, cobalt-free or cobalt-reduced NMC (LiNi$_x$Mn$_y$Co$_z$O$_2$) formulations with high nickel content (NMC811, NMC90) are being pursued to reduce cost and supply risk while improving energy density. LiFePO$_4$ (LFP) has re-emerged strongly as a safe, low-cost, cobalt-free cathode with excellent cycle life, particularly for stationary storage and lower-range EVs.

\subsection{Manufacturing and Component Challenges}

Component-level challenges include improving electrode adhesion, electronic conductivity, and electrolyte compatibility. The replacement of PVDF/NMP binder systems with water-based alternatives (biopolymer gums such as xanthan or carboxymethylcellulose) is an active area, motivated both by environmental concerns and by the need to reduce production costs and eliminate solvent recovery steps.

At the battery assembly level, the absence of any metallic contamination is a paramount constraint, since even sub-millimetre metal particles can provoke internal short circuits leading to thermal runaway. Clean-room manufacturing environments, continuous monitoring, and automated optical inspection are indispensable.

\subsection{Cost Trajectory}

The price of Li-ion battery packs has fallen dramatically: from approximately \$780/kWh in 2013 to around \$139/kWh in 2023. The widely cited cost target of \$100/kWh is considered the threshold at which electric vehicles reach purchase price parity with internal combustion vehicles without subsidies. Continued cost reduction requires simultaneously improving energy density, reducing reliance on expensive materials, and scaling manufacturing.

\subsection{Safety, Recycling, and End-of-Life}

Safety remains a concern due to the flammability of liquid electrolytes and the risk of thermal runaway. Solid-state electrolytes — ceramics, sulfides, or polymer-ceramic composites — promise intrinsically safer batteries by eliminating flammable liquids and enabling lithium metal anodes, but achieving sufficient ionic conductivity and stable electrode/electrolyte interfaces at room temperature remains an unsolved challenge.

Recycling is increasingly critical both for resource recovery and for reducing the environmental footprint of battery production. Current recycling routes (pyrometallurgy, hydrometallurgy, direct recycling) recover metals to varying degrees but remain costly and energy-intensive. The EU Battery Regulation introduces mandatory recycled content targets and collection obligations.

\section{Future Perspectives}

\subsection{Beyond Li-ion: Sodium-Ion Batteries}

Sodium-ion (Na-ion) batteries follow the same intercalation principle as Li-ion but replace Li$^+$ with Na$^+$. Sodium is vastly more abundant and geographically distributed than lithium, and Na-ion cells can be assembled without copper current collectors (since sodium does not alloy with aluminium at low potentials, aluminium can be used for both electrodes). Water-based electrode manufacturing is also feasible. Current Na-ion cells reach around 100\,Wh/kg with cycle lives exceeding 2000 cycles — lower than state-of-the-art Li-ion but improving rapidly. Cathode materials include sodium vanadium fluorophosphates (NaVPO$_4$F and variants); hard carbon is the preferred anode. Industry is moving fast: BYD announced in February 2026 a third-generation Na-ion pack targeting 10,000 cycles, alongside a solid-state battery programme for 2027 EVs.

\subsection{Solid-State Batteries}

Solid-state batteries replace the liquid electrolyte with a solid ionic conductor (oxide, sulfide, or polymer-based). They promise higher energy density (by enabling lithium metal anodes), intrinsic safety (no flammable liquid), and potentially wider operating temperature ranges. Key unresolved challenges include: achieving room-temperature ionic conductivity comparable to liquid electrolytes ($\sim 10^{-2}\,\text{S/cm}$), ensuring mechanically stable and low-resistance electrode/electrolyte interfaces over thousands of cycles, and scaling up manufacturing.

\subsection{A Broader Perspective}

Battery technology is ultimately only one part of the energy transition. Research cited in the presentation (Barrett et al., \textit{Nature Energy}, 2022) estimates that current per-capita energy consumption in developed countries averages $\sim$6\,kW/person, while sustainability analyses suggest a target of $\sim$2\,kW/person is necessary. Improving energy storage is therefore inseparable from the broader imperative to reduce absolute energy demand.




\section{Podcast -- Professeure Nathalie Job, Université de Liège}

\subsection{Quel est le rôle du carbone dans la vie de recherche de Nathalie Job ?}

Le carbone est un matériau aux applications très nombreuses : graphite, graphène, diamant, fibres de carbone, charbons actifs pour la purification de l'eau, etc. Nathalie Job a commencé par travailler sur des matériaux carbonés poreux appelés \textit{xérogels de carbone}, utilisés en catalyse. Le carbone est aussi conducteur électrique, ce qui le rend précieux pour la fabrication d'électrodes en électrochimie. C'est donc par le carbone qu'elle est entrée dans la recherche, et ce matériau constitue le fil conducteur de l'ensemble de ses activités.

\subsection{Qu'est-ce qu'un xérogel ?}

Le terme \textit{xérogel} vient du grec \textit{xéros} signifiant « sec » : c'est donc un gel séché. On part de produits en solution qui gélifient, puis on les sèche et on les pyrolyse, ce qui produit une structure carbonée tridimensionnelle et poreuse. Cette structure est similaire au charbon actif utilisé dans les filtres d'aquarium ou les systèmes d'épuration d'eau, à la différence qu'on y dépose des métaux pour réaliser de la catalyse en phase gazeuse ou liquide.

\subsection{Quel était l'enjeu des xérogels en termes de surface et de porosité ?}

L'enjeu principal n'était pas tant d'augmenter la surface spécifique (le charbon actif en possède déjà une très grande), mais de \textit{contrôler la texture poreuse} du matériau. Pour certaines applications de catalyse en phase gazeuse, il est nécessaire de disposer de méso- ou macropores bien structurés (entre 10 et 100~nm) afin d'améliorer le transfert de matière vers les sites actifs à l'intérieur du catalyseur. En maîtrisant cette texture, on obtient de meilleures activités catalytiques.

\subsection{Comment Nathalie Job est-elle entrée dans le domaine des piles à combustible ?}

Le carbone étant à la fois riche chimiquement (on peut y déposer des métaux, de l'azote, de l'oxygène) et conducteur électrique, il entre naturellement dans la fabrication des électrodes. Les piles à combustible à hydrogène basse température utilisent des catalyseurs constitués de métaux précieux dispersés sur du noir de carbone. Nathalie Job a transposé ses travaux sur les xérogels de carbone nanostructurés à ce domaine, en cherchant à améliorer le transfert de matière dans les électrodes de ces piles.

\subsection{Comment fonctionne une pile à combustible ?}

Une pile à combustible sépare deux demi-réactions électrochimiques, comme dans une pile de Daniell :
\begin{itemize}
    \item À l'anode : oxydation de l'hydrogène en protons.
    \item À la cathode : réduction de l'oxygène en eau.
\end{itemize}
Les protons transitent de l'anode vers la cathode à travers un électrolyte solide polymère (membrane), tandis que les électrons circulent dans le circuit extérieur, générant ainsi un courant électrique. Pour fonctionner, le système doit gérer simultanément le flux de gaz, l'eau produite, les électrons (conduits par le carbone solide) et les protons (conduits par la membrane). La structure poreuse des électrodes est donc cruciale.

\subsection{Pourquoi les piles à combustible sont-elles si coûteuses ?}

Trois facteurs principaux expliquent ce coût élevé :
\begin{enumerate}
    \item Le catalyseur : le platine (souvent allié au cobalt ou au nickel) est indispensable et très cher.
    \item L'électrolyte : la membrane polymère utilisée, le Nafion, reste très onéreuse.
    \item L'absence d'économies d'échelle : la production en série n'a pas encore atteint un volume suffisant pour réduire significativement les coûts de fabrication.
\end{enumerate}

\subsection{Quels matériaux à longue durée de vie cherche-t-on dans les piles à combustible ?}

Deux problèmes principaux se posent au niveau du catalyseur. D'abord, \textit{l'activité} : on cherche des catalyseurs plus actifs afin de réduire les pertes de tension (et donc d'énergie) à l'électrode, voire de trouver des alternatives moins coûteuses au platine. Ensuite, \textit{la durabilité} : le couple platine-carbone n'est pas stable thermodynamiquement sous tension et en présence d'hydrogène et d'oxygène. Les électrodes ne durent pas assez longtemps pour couvrir toute la durée de vie d'une automobile, ce qui implique un remplacement coûteux de la pile en cours de vie du véhicule.

\subsection{Quelle est la chaîne de rendement dans une pile à combustible ?}

Les piles à combustible atteignent des rendements de 60 à 65\% en conditions réelles, contre un maximum théorique (à courant nul) d'environ 95\%. Lorsqu'on tire du courant, plusieurs types de pertes apparaissent :
\begin{itemize}
    \item Les \textit{pertes d'activation} : liées à l'énergie d'activation de la réaction catalytique, surtout à faible courant.
    \item Les \textit{pertes ohmiques} : dues à la résistance électrique des composants (loi d'Ohm).
    \item Les \textit{pertes de transport de matière} : à courant élevé, les réactifs peinent à atteindre les sites actifs. C'est ici qu'intervient l'intérêt des xérogels de carbone nanostructurés.
\end{itemize}

\subsection{Que fait-on concrètement dans le projet Be Hyfe pour prolonger la durée de vie des catalyseurs ?}

Deux axes sont poursuivis. Le premier consiste à améliorer les \textit{propriétés de surface du carbone} : les noirs de carbone industriels sont très désordonnés et donc sensibles à la corrosion. L'idée est de déposer une seconde couche graphitisée sur les xérogels de carbone poreux (par dépôt chimique en phase vapeur, par exemple) afin de combiner cristallinité et nanostructuration poreuse. Le second axe porte sur la \textit{phase active} : au lieu des nanoparticules de platine classiques (3--4~nm, très mobiles), on développe des particules creuses de platine-cobalt de 30 à 40~nm de diamètre, beaucoup plus stables car moins sujettes à la coalescence ou à la redissolution.

\subsection{L'objectif est-il de limiter la masse de platine utilisée ?}

L'objectif premier est de pouvoir utiliser le platine sur la plus longue durée possible, c'est-à-dire d'améliorer la durabilité du système. Avoir une grande quantité de platine dans un système durable n'est pas un problème en soi. En revanche, avoir beaucoup de platine dans un système qui vieillit rapidement est très problématique. En parallèle, allier le platine à d'autres éléments permet également d'en réduire la quantité pour un rendement équivalent.

\subsection{Existe-t-il une alternative viable au platine ?}

À l'échelle industrielle, non. Des recherches actives portent sur des catalyseurs de type fer-carbone-azote, qui fonctionnent à petite échelle, mais leur maturité technologique est insuffisante : la fabrication d'électrodes complètes est difficile et leur stabilité reste inférieure à celle du platine sur carbone. C'est néanmoins un vaste champ de recherche.

\subsection{Quel est le point de vue de Nathalie Job sur l'utilisation de l'hydrogène ?}

Il faut considérer l'ensemble de la chaîne. La production d'hydrogène vert par électrolyse à partir d'énergies renouvelables, suivie d'une conversion en électricité via une pile à combustible, conduit à un rendement global d'environ 25\%. Cela n'est pas pertinent pour toutes les applications. Pour les petites voitures urbaines, la batterie est plus efficace. En revanche, pour des applications à forte puissance (camions, trains, bateaux) ou pour le stockage stationnaire de longue durée, l'hydrogène est plus approprié. Il ne faut pas généraliser et prétendre que l'hydrogène remplacera le pétrole dans tous les usages : les besoins en électricité pour produire suffisamment d'hydrogène seraient considérables.

\subsection{Comment fonctionne l'électrolyse ? Est-ce l'inverse de la pile à combustible ?}

Oui, c'est exactement l'inverse : au lieu de recombiner l'hydrogène et l'oxygène pour produire de l'électricité, on injecte de l'électricité pour dissocier l'eau en hydrogène et oxygène. Il existe des électrolyseurs PEM (à membrane échangeuse de protons) et des électrolyseurs alcalins, analogues inverses de leurs piles respectives. Toutefois, les matériaux optimaux ne sont pas exactement les mêmes dans les deux sens car les pertes ne sont pas symétriques entre oxydation et réduction, ce qui rend difficile la conception de systèmes réversibles efficaces.

\subsection{Dans quelles applications un système réversible électrolyse/pile à combustible serait-il pertinent ?}

La principale motivation est la \textit{compacité} : sur un bateau, par exemple, il serait avantageux d'avoir un unique système capable à la fois de stocker l'électricité sous forme d'hydrogène et de la restituer via une pile à combustible, plutôt que deux systèmes distincts. Cependant, atteindre une efficacité comparable dans les deux sens reste un défi non résolu à ce jour.

\subsection{Quel est le lien entre les activités sur les batteries et sur les piles à combustible ?}

Le lien est le carbone. Les batteries lithium-ion commerciales contiennent du graphite à l'anode, et les problématiques de structuration des électrodes sont similaires dans les deux domaines : mélanger des matériaux actifs, un électrolyte et un liant, puis les déposer sur un collecteur de courant. Nathalie Job est entrée dans les batteries via les matériaux nanostructurés (gels de carbone), en étudiant l'impact de la texture nanométrique sur la fabrication et les performances des électrodes.

\subsection{Comment fonctionne une batterie lithium-ion ?}

Une batterie lithium-ion exploite la capacité des ions lithium (Li$^+$) à s'intercaler à l'intérieur de matériaux solides. À la charge, les ions lithium sont extraits de la cathode et insérés dans l'anode (graphite) ; les électrons transitent simultanément par le circuit extérieur. À la décharge, le processus s'inverse spontanément : les ions retournent à la cathode et les électrons circulent dans le circuit, fournissant le courant. L'innovation fondamentale, récompensée par un prix Nobel de chimie en 2019, réside dans ces \textit{matériaux d'intercalation} qui permettent de stocker les charges dans toute la masse du matériau, et non seulement en surface.

\subsection{Comment sait-on qu'une batterie est complètement chargée ?}

Il n'existe pas de mesure directe de l'état de charge. On ne peut mesurer que la tension et le courant. L'état de charge est estimé par le \textit{Battery Management System} (BMS), qui s'appuie sur des données issues de cycles de charge/décharge préétablis par le fabricant. Le lien entre tension et état de charge n'est pas direct : pour certains matériaux présentant un plateau de tension très plat, la corrélation est quasi nulle, ce qui rend le contrôle encore plus complexe.

\subsection{Quels sont les effets de la température sur les batteries ?}

Une baisse de température réduit la conductivité ionique de l'électrolyte, ce qui peut amener le BMS à interrompre la charge ou la décharge plus tôt, réduisant ainsi l'autonomie apparente. La batterie répond également plus lentement. À l'inverse, une température trop élevée pose des problèmes de sécurité : l'électrolyte (composé de solvants organiques très inflammables) peut se dégrader en produisant des gaz (CO$_2$, H$_2$), entraîner un gonflement de la cellule, une rupture de l'enveloppe, et potentiellement un incendie au contact de l'air.

\subsection{Existe-t-il une alternative aux électrolytes liquides inflammables ?}

Oui, c'est l'objet des \textit{batteries tout-solides} : on remplace l'électrolyte liquide par un polymère solide ou une céramique, qui ne peuvent pas brûler. Cependant, la conductivité ionique des électrolytes solides est actuellement bien inférieure à celle des liquides, ce qui entraîne une perte de performance incompatible, pour l'instant, avec les exigences de l'automobile. Les risques d'incendie des batteries actuelles restent donc non négligeables.

\subsection{Quel est le point de vue de Nathalie Job sur la sécurité des batteries domestiques ?}

Elle déconseille fortement d'installer chez soi des batteries usagées (notamment issues de trottinettes) dont l'état de dégradation est inconnu. Les incendies de batteries sont difficiles à éteindre (l'eau aggrave la situation), et les entreprises de recyclage de vélos électriques font face à ce problème en pratique, notamment aux Pays-Bas. La législation sur la sécurité des batteries domestiques devra évoluer significativement.

\subsection{Pourquoi les batteries perdent-elles de leurs performances avec le temps ?}

Plusieurs mécanismes sont en jeu. D'abord, les matériaux subissent des \textit{sollicitations mécaniques} répétées : l'intercalation/désinsertion du lithium fait gonfler et rétrécir les particules cycliquement, ce qui peut les fissurer et les désolidariser. Ensuite, si l'on extrait trop de lithium de certains matériaux de cathode (comme le LCO), leur structure cristalline change de manière irréversible. Le graphite s'exfolie progressivement. Enfin, l'électrolyte se dégrade aux interfaces, produisant des gaz indésirables. Le BMS tient compte de ce vieillissement en réduisant progressivement la plage de charge/décharge utilisée.

\subsection{La plage de charge et décharge est-elle réduite avec le temps ?}

Oui, c'est exactement cela. Ce que l'on perçoit comme « 0 à 100\% » sur un appareil correspond en réalité à une plage sécurisée définie par le BMS, et non aux limites physiques absolues du matériau. Avec le vieillissement, cette plage se réduit encore davantage, ce qui se traduit par une perte d'autonomie perceptible pour l'utilisateur.

\subsection{Faut-il éviter de charger et décharger trop rapidement ?}

Oui. Charger ou décharger très rapidement impose des contraintes plus importantes aux matériaux et accélère leur vieillissement. Solliciter les batteries à des tensions extrêmes (charge complète prolongée) augmente certes l'autonomie momentanée, mais dégrade le matériau plus vite. Idéalement, il faudrait éviter les charges et décharges trop profondes et trop rapides.

\subsection{Que fait-on pour optimiser le BMS et faire le lien avec la charge rapide ?}

En collaboration avec l'équipe de Michel Kinnaert au laboratoire SAAS de l'ULB, les travaux portent sur la modélisation mathématique du comportement de la batterie, alimentée par des données expérimentales fournies par le laboratoire de Nathalie Job. L'objectif est de développer des protocoles de charge plus intelligents : par exemple, varier le courant de façon optimale pour raccourcir la durée de charge sans endommager les matériaux, au lieu du simple protocole courant constant / tension constante. La \textit{charge rapide} ne correspond en pratique qu'à la première phase (courant constant), sans aller jusqu'à la saturation — ce qui préserve la durée de vie mais réduit l'autonomie maximale.

\subsection{Quelles autres avancées perçoit-on pour les batteries dans le futur ?}

Les évolutions seront dictées par la nécessité, notamment la raréfaction et le coût croissant des matériaux critiques (lithium, cobalt, cuivre, graphite). Les pistes incluent : le remplacement du LCO par des phosphates de fer lithium (LFP) ou des NMC (nickel-manganèse-cobalt), le remplacement du graphite par des oxydes de titane, le remplacement du cuivre par du carbone, l'élimination du liant PVDF (fluoré, difficile à recycler et nécessitant des solvants organiques nocifs), ainsi que le développement d'électrolytes solides. Le recyclage des matériaux est également un enjeu majeur.

\subsection{Est-il réaliste de ne plus vendre de voitures à combustion à partir de 2035 ?}

Nathalie Job se montre sceptique. Remplacer les moteurs à combustion par des véhicules électriques à batterie ne fait que substituer une dépendance (pétrole) par une autre (cobalt, lithium, cuivre, fabrication asiatique). La chaîne de fabrication des batteries n'est pas aussi écologique qu'on le prétend souvent. Il ne faut pas croire aux solutions uniques : chaque application a ses propres contraintes, et une transition aussi radicale imposée par décret sans changement de comportement ne suffira pas à résoudre le problème environnemental.

\subsection{Le projet de « peindre » des batteries : de quoi s'agit-il ?}

Dans le cadre du projet Batwal, la contribution du laboratoire était de développer des \textit{électrodes fabriquées par spray coating à l'eau}. Le principe : pulvériser une suspension aqueuse de matériaux actifs (LFP en cathode, Li$_4$Ti$_5$O$_{12}$ en anode) sur un collecteur de courant conducteur, en remplaçant le liant PVDF par de la gomme de xanthane (un polysaccharide alimentaire) et le solvant organique par de l'eau. Les performances obtenues sont équivalentes aux électrodes conventionnelles, tout en éliminant les solvants nocifs et en facilitant grandement le recyclage.

\subsection{Une spin-off a-t-elle été lancée suite au dépôt de brevet ?}

Oui. Un projet FESO financé par la Région wallonne est en cours, visant à valoriser les technologies d'électrodes développées au laboratoire. La spin-off \textit{Green Li-ion} est en cours de création. Elle s'appuie sur deux technologies brevetées : le spray coating à l'eau avec gomme de xanthane, et un procédé de spray coating du LCO sans aucun additif (ni PVDF, ni carbone conducteur), permettant d'obtenir une électrode pure, facilement séparable et recyclable.

\subsection{Quels matériaux sont particulièrement sensibles et font l'objet de recherches ?}

Les matériaux NMC (oxydes de nickel, manganèse et cobalt), qui remplacent progressivement le LCO dans l'automobile pour réduire la teneur en cobalt, sont difficiles à traiter avec des procédés aqueux car ils se dégradent au contact de l'eau. Les matériaux conducteurs carbonés (noir de carbone, nanotubes) s'agglomèrent également dans l'eau. Des recherches sont en cours pour modifier la chimie de surface de ces matériaux afin de les rendre compatibles avec des procédés aqueux.

\subsection{L'industrialisation du procédé a-t-elle été envisagée ?}

Oui, la faisabilité industrielle est gardée en tête dès la conception des procédés. Les travaux au laboratoire ont notamment porté sur la simplification maximale des étapes de synthèse et sur la mise au point de procédés continus (par opposition aux batchs). Un procédé de fabrication en continu de billes de gel de carbone a par exemple été développé, démontrant la faisabilité de l'industrialisation. Le choix de matières premières non critiques et de procédés simples est systématiquement privilégié.

\subsection{La technologie seule peut-elle résoudre le problème énergétique ?}

Non. Nathalie Job est convaincue que la technologie seule ne suffira pas. Les courbes historiques de consommation énergétique montrent qu'aucune source n'a jamais été abandonnée : on ne fait qu'en ajouter de nouvelles. Si la technologie ne s'accompagne pas de changements de comportement, elle ne fera qu'alimenter une consommation croissante. La sobriété énergétique est indispensable et devrait être anticipée plutôt que subie.

\subsection{Quels changements de croyances ou d'habitudes Nathalie Job a-t-elle entrepris ces dernières années ?}

Elle a dû apprendre à \textit{déléguer} et à lâcher prise sur le contrôle direct des expériences, en passant du rôle de chercheuse exécutante à celui de directrice d'équipe. Cette transition — superviser des doctorants, gérer des projets et des cours tout en abandonnant la paillasse — n'a pas été aisée, mais s'avère inévitable pour mener de front toutes les responsabilités d'un poste de professeure ordinaire.

\subsection{Quels livres conseille Nathalie Job ?}

Elle recommande \textit{La Guerre éternelle} de Joe Haldeman, un roman de science-fiction des années 1970. L'humanité y entre en guerre contre une civilisation extraterrestre jamais rencontrée, et les affrontements sont conditionnés par des sauts dans l'espace-temps, si bien qu'aucun des deux camps ne sait d'avance lequel est technologiquement supérieur. C'est un roman anti-guerre, moins connu que d'autres classiques du genre, mais remarquable.

\subsection{Quelle phrase Nathalie Job écrirait-elle sur tous les tableaux de l'université ?}

« \textbf{Votre consommation énergétique : 6~kW.} » Cette valeur représente la puissance moyenne consommée en continu par un habitant des pays développés, en tenant compte de tous les usages (transport, alimentation, chauffage, vêtements, etc.). Ce chiffre illustre l'impossibilité de couvrir l'ensemble des besoins énergétiques actuels uniquement avec des énergies renouvelables intermittentes, et plaide pour une réduction sérieuse de la consommation.

\subsection{Quel hobby a influencé la personnalité de Nathalie Job ?}

La \textit{voile en compétition}, pratiquée à un niveau élevé avec son père, notamment en mer du Nord. La leçon principale : ne jamais abandonner et toujours terminer ce qu'on commence, même en cas d'échec momentané. Cette philosophie — finir la course quelles que soient les conditions, sauf blessure ou avarie — a forgé son caractère de chercheuse : un projet refusé se réécrit, un article rejeté s'améliore et se resoumet.